\section{Multi-resolution field synthesis}
\label{sec:multires}

The expensive step in a monolithic implementation is the global $N^3$ FFT
(with $N=6144$, a single Fourier-space field is 1.7~TB in double-precision
complex), repeated for the density, the filter bank, and seven 2LPT component
fields. \ppjl\ replaces the global transform with a two-level decomposition
whose key property is \emph{exactness of the noise}: the union of tiles
realizes the same white-noise field as the monolithic grid, independent of the
tiling, so tile decomposition is a numerical method, not a change of
realization.

\subsection{Tiling-independent noise}
\label{sec:multires-noise}

The $z=0$ white noise is defined by a counter-based (Threefry) random number
generator \citep{Salmon2011} keyed on the \emph{global lattice index}: cell
$(i,j,k)$ of the $N^3$ grid receives a unit Gaussian computed directly from
its flattened index and the run seed. Any subvolume can therefore generate its
own noise values without communication, and the same seed yields the same
global field for every tiling, every device count, and both the halo and
map-making passes (which exploit this determinism to re-synthesize identical
fields in a separate pass, Section~\ref{sec:maps}).

\subsection{Coarse field and tile residuals}
\label{sec:multires-split}

The box is divided into ${n_{\rm tile}}^3$ tiles of core size
$n_{\rm sub}^3$, each extended by a buffer of $n_{\rm buff}$ cells on every
face ($n_{\rm mesh}=n_{\rm sub}+2n_{\rm buff}$); buffers absorb the finite
range of the transfer function and of halo exclusion across tile faces, and
peaks are only accepted in the core. Long-wavelength modes are carried by a
global coarse grid of size $M^3$ with $M = n_{\rm tile}\times f_c$ for a
coarsening factor $f_c$:
\begin{enumerate}
  \item the coarse noise is the block average of the fine noise over
        $(N/M)^3$ blocks (normalized to unit variance), computed directly
        from the Threefry stream;
  \item the coarse field is synthesized with a standard periodic FFT
        convolution with $\sqrt{P(k)}$ on the $M^3$ grid, then interpolated
        onto each tile's fine lattice;
  \item each tile computes its \emph{residual} noise --- fine noise minus the
        (nearest-neighbour upsampled) coarse block mean --- which has exactly
        zero mean in every coarse cell. This is a Hoffman--Ribak-style
        constrained split: because the residual carries no power at
        wavelengths longer than the coarse cell, its convolution with
        $\sqrt{P(k)}$ is short-ranged and can be evaluated with an
        \emph{isolated} (zero-padded, non-periodic) FFT over the tile plus a
        shell of surrounding cells, then trimmed to the tile.
\end{enumerate}
The tile's total field is the interpolated coarse field plus the isolated
residual convolution. The same decomposition is applied to every derived
field (smoothed densities, displacement and 2LPT component fields), so each
tile runs the entire pipeline of Section~\ref{sec:algorithm} independently.
\todo{Fig.~F1: tiling/buffer schematic; Fig.~F2: splice cross-section +
power-spectrum continuity across the coarse--fine transition}

\subsection{A resolution criterion for velocity-weighted products}
\label{sec:multires-velocity}

The split is benign for density-derived statistics at any practical $f_c$:
mis-assigned power near the coarse--fine transition band moves structure
coherently within a tile without changing local densities, and all
catalogue-level and lensing statistics in
Sections~\ref{sec:valcat}--\ref{sec:valmaps} are insensitive to $f_c$.
Velocity-weighted observables are different. The velocity field weights the
density modes by $\propto k^{-1}$, so the transition band contributes a much
larger share of the signal, and any decoherence between the coarse
interpolation and the tile residuals in that band appears directly in
line-of-sight momentum statistics.

We measured this with the kSZ field maps of Section~\ref{sec:maps}: at our
production geometry ($n_{\rm tile}=16$), a coarse grid of $64^3$
($f_c=4$, Nyquist $k_{\rm Nyq}=0.038\,h\,{\rm Mpc}^{-1}$) produces a
2--3$\times$ excess in the kSZ angular power spectrum at $\ell\lesssim 350$
--- exactly the multipoles that map to the transition band at the
$\tau$-weighted distances --- while leaving the lensing convergence control
map unchanged at the few-per-cent level. Increasing the coarse grid to
$256^3$ ($f_c=16$, $k_{\rm Nyq}=0.154$) reduces the excess to
$\lesssim$40\%, and $512^3$ ($f_c=32$, $k_{\rm Nyq}=0.31$) converges the kSZ
spectrum, changing the density control by $\lesssim$few per cent
(Section~\ref{sec:valmaps}; Fig.~\todo{F8}). We therefore adopt the criterion
\begin{equation}
  k_{\rm Nyq}^{\rm coarse} \;=\; \frac{\pi M}{L_{\rm box}}
  \;\gtrsim\; 0.3\,h\,{\rm Mpc}^{-1}
  \label{eq:cfrule}
\end{equation}
for any velocity- or potential-weighted product (kSZ, ISW, moving lens),
while $f_c=4$ suffices for catalogues and density-derived maps. The coarse
FFT remains negligible in cost at $512^3$, so the criterion carries no
performance penalty; it is a correctness requirement that we expect applies
to any tile-based approximate method that synthesizes velocities from a
decomposed displacement field.
